This report has examined \emph{tabular data understanding} with large language models through two complementary pillars (schema understanding and data understanding) and situated our own contribution within them.

Chapter~\ref{ch:ch1} framed the field around these two pillars. It began from the representation mismatch that makes tabular reasoning hard: because language models consume one-dimensional token sequences while tables are inherently two-dimensional, every pipeline must commit to a representation, and this choice (serialization, data schema, image, a specialized table encoder, or a graph) conditions everything downstream. The survey organized the field into schema understanding, the mapping of natural-language intent onto a structural blueprint, and data understanding, the problem of reasoning faithfully over table content, treating representation within the data pillar and singling out graph representation as a solid but so-far small-scale and document-centric approach whose extension to large tables is a central opportunity. Across both, a consistent picture emerged: single-pass methods have given way to multi-step pipelines that decompose, retrieve, verify, and correct, and the binding constraints are increasingly architectural (how structure and evidence are delivered to the model) rather than a matter of raw model capacity.

Chapter~\ref{ch:ch3} presented our contribution to the schema-understanding pillar, Reflect-SQL~\cite{jain2027reflectsql}, which threads self-reflection through all three stages of a Text-to-SQL pipeline: it constructs a semantic knowledge base over the schema, retrieves relevant elements iteratively, and synthesizes queries under syntactic and multi-dimensional semantic validation, using an LLM-as-a-judge to score correctness without a gold reference. Grounding reflection in the checkable signals that the SQL setting affords (execution, syntax, and result--intent entailment) allows the system to reach $72.03\%$ execution accuracy on the industry-scale BIRD benchmark, and demonstrates that reference-free, self-critical generation is a viable path for structured tasks.

Chapter~\ref{ch:future} then set out the research agenda that follows. On the schema side, the priorities are to make the knowledge base continually evolving rather than static, to scale retrieval by addressing pre-computed structure rather than searching at query time, to extend faithful synthesis to cross-database joins, and to make reflection both reliable and inexpensive. On the data-understanding side, our next pillar, the target task is Table Question Answering and the lead direction is to represent large tabular data at scale, concretely as a graph, so that evidence retrieval becomes addressing rather than search, with robustness on complex table structures and higher-order reasoning as the targets and a large-scale Table QA benchmark as the instrument to measure them.

The unifying thesis of this report is that progress in tabular data understanding depends less on ever-larger models than on the architecture that supplies them with the right structure and the right evidence, and that this is above all what will let the field cross from small, curated benchmarks to the large, obscure, obsolescent data of real enterprises. Reflect-SQL is a concrete step in that direction for schema understanding; carrying the same principle into data understanding, so that both pillars hold at industry scale, is the trajectory we intend to pursue.
